\section{Coefficient dictionaries and passage domains}
\label{app:coefficient-dictionaries}

We express each coefficient limit in a chart where the corresponding
passage estimate is uniform. The substitutions are exact, so the estimates
apply to the original coefficient fibres. We then choose overlapping
chart ranges in the order required by Section~\ref{sec:coefficient-limits}.
All constants precede the common raw cap and final coefficient
restrictions in Lemma~\ref{lem:common-raw-receiver}.

\subsection{Proof of Lemma~\ref{deg:lem-small-source-routing}}
\label{app:proof-small-source-routing}

\begin{proof}
If \(p\le0\), use Lemma~\ref{deg:lem-raw-jacobian}. For \(p>0\),
first fix the positive-quadratic ghost and root-free reserves, and then
the independent far-root rectangles.

Fix a small number \(\mathfrak q_{\rm far}>0\) for which
Proposition~\ref{deg:prop-positive-far} applies.  On the negative
\(\gamma\)-side, Proposition~\ref{deg:prop-degree-infinity} and the
dictionary below provide numbers \(\mathfrak q_- >0\) and
\(\Pi_->1\).  Put
\begin{equation}
 \mathfrak q_\flat=\frac12\min(\mathfrak q_{\rm far},\mathfrak q_-),\qquad
 \Pi_{\rm far}=\max\{\lambda_0^{-1/2},\Pi_-,2\sqrt{12}\},
 \qquad \Pi_{\max}=2\Pi_{\rm far},
 \label{deg:small-constant-1}
\end{equation}
where \(\lambda_0\) is the constant in
\eqref{deg:positive-far-range}. Apply the whole-carrier estimate of
Corollary~\ref{deg:cor-positive-quadratic-whole-carrier} and
Proposition~\ref{deg:prop-negative-gamma} on
\(0<\pi\le\Pi_{\max}\).

For \(\mathfrak q=k\pi\ge \mathfrak q_\flat\), the degree-loss dictionary given below is
uniform once \(k\le k_d(\mathfrak q_\flat)\).  Choose
\(k_\flat>0\) so that \(2k_\flat\) is smaller than this number and than
the \(k\)-bounds in both compact-\(\pi\) carriers, including the
whole positive-carrier bound \(k_0\le k_{\rm tail}\).  Put
\begin{equation}
 G_{\max}=k_\flat^{-2/3}.
 \label{deg:small-Gmax}
\end{equation}
For this fixed \(G_{\max}\), choose \(U_d\) so that the large-\(U\)
dictionary below belongs to Proposition~\ref{deg:prop-degree-drop}, and
then choose \(U_{\rm cub}>U_d+1\).  Apply Proposition~\ref{deg:prop-compact-cubic} on
\(0\le U\le U_{\rm cub}\), \(|G|\le G_{\max}\). The bounded
positive-quadratic \(k\to0\) face remains in the first compactification;
it is distinct from this compact cubic rectangle.

We first check that the two degree-loss substitutions reproduce the
source coefficients exactly.  In the \(\gamma\)-dominant
chart \eqref{deg:gamma-chart}, let
\(\epsilon=\operatorname {sgn}\gamma\) and put \(t=k\xi\).
The exact polynomial is
\begin{equation}
 F_{\epsilon,\mathfrak q,k}(t)=t^3+3\epsilon t^2-\mathfrak q t+k^2.
 \label{deg:gamma-degree-polynomial}
\end{equation}
For \(\mathfrak q\ge \mathfrak q_\flat\) and \(k\) small it has roots
\(-r<t_a<t_c\), all three simple.  Vieta's formulas are
\[
 t_a+t_c=r-3\epsilon,\qquad
 rt_at_c=k^2,\qquad
 \mathfrak q=r(r-3\epsilon)-\frac{k^2}{r}.
\]
Define
\begin{equation}
 \begin{gathered}
 \tau=p^2,\qquad
 \kappa=\frac{k^2}{\mathfrak q^{3/2}},\qquad
 \mathcal R=\frac{r}{\sqrt{\mathfrak q}},\qquad
 B=\frac{\mathfrak q}{r^2},\qquad \omega=\frac{\sqrt{\mathfrak q}}{r},\\
 \nu=\frac{\epsilon}{r},\qquad
 a=\frac{t_a}{r},\qquad c=\frac{t_c}{r},\qquad
 A_d=\frac{B^{3/2}}{c}.
 \end{gathered}
 \label{deg:gamma-degree-dictionary}
\end{equation}
These identities imply \eqref{deg:degree-graph} exactly.  Moreover
\[
 (\Lambda_1,\Lambda_2,\Lambda_3)
 =(1,\kappa^2,\epsilon/\pi^2),
\]
so
\[
 (\sqrt\tau\Lambda_1,\tau\Lambda_3,\tau^{3/2}\Lambda_2)
 =(p,\gamma,s).
\]
For \(\mathfrak q\in[\mathfrak q_\flat,\infty]\), all positive-unit quantities in
Proposition~\ref{deg:prop-degree-drop} have compact closure:
\(r/\sqrt \mathfrak q\to1\) and
\((B,\omega,c,\nu,A_d)\to(1,1,1,0,1)\) as \(\mathfrak q\to\infty\).

On the negative-\(\gamma\) far face \(\mathfrak q\downarrow0\),
\(\pi\to\infty\), use the roots of
\[
 F_{-,\mathfrak q,k}(t)=t^3-3t^2-\mathfrak q t+k^2.
\]
The exact regular coordinates of Proposition
\ref{deg:prop-degree-infinity} are
\begin{equation}
 \begin{gathered}
 a=\frac{t_a}{r},\qquad j=\frac{r}{t_c},\qquad
 \kappa=a,\qquad \mathcal R=j,\qquad d=B=A_d=1,\\
 \nu=-\frac1r,\qquad \omega=1-a+ja,\qquad
 \tau=\frac{gt_c}{a}.
 \end{gathered}
 \label{deg:negative-far-dictionary}
\end{equation}
Vieta's formulas give
\[
 t_a+t_c=r+3,\qquad rt_at_c=k^2,\qquad
 \mathfrak q=r(r+3)-\frac{k^2}{r},\qquad
 \pi^2=\frac{t_c}{a}(1-a+ja)^2,\qquad
 \tau=\frac{p^2}{\omega^2}.
\]
Hence
\[
 r\sim \mathfrak q/3,\quad t_c\to3,\quad a\sim3/\pi^2,\quad
 j\sim \mathfrak q/9,\quad \tau\sim p^2.
\]
Thus \((j,\kappa,\tau)\to0\), the positive-unit entries stay in
compact ranges, and \(j\nu=-1/t_c=(j-ja-1)/3\) stays bounded.
The same source coefficients are recovered exactly. Fix
\(\mathfrak q_-\), then \(\Pi_-\), so that
\(\omega\ge\omega_{\min}>0\), and let \(\tau_-\) be the
reciprocal radial bound. The rectangle
\begin{equation}
 0<\mathfrak q\le \mathfrak q_-,\qquad \pi\ge\Pi_-,\qquad
 p^2\le\tau_-\omega_{\min}^2
 \label{deg:negative-far-rectangle}
\end{equation}
lies within the range of Proposition~\ref{deg:prop-degree-infinity}.
The estimate therefore holds throughout this rectangle, independently
of the relative rates of its parameters.

For the complementary cubic chart \eqref{deg:cubic-chart}, the large
\(U\) polynomial
\[
 y^3+3Gy^2-Uy+1
\]
has three simple roots \(-r<t_a<t_c\) when
\(U\ge U_d(G_{\max})\).  Set
\begin{equation}
 \begin{gathered}
 \tau=p^2,\qquad \kappa=U^{-3/2},\qquad
 \mathcal R=\frac{r}{\sqrt U},\qquad
 B=\frac{U}{r^2},\qquad \omega=\frac{\sqrt U}{r},\\
 \nu=\frac{G}{r},\qquad a=\frac{t_a}{r},\qquad
 c=\frac{t_c}{r},\qquad A_d=\frac{B^{3/2}}{c}.
 \end{gathered}
 \label{deg:cubic-degree-dictionary}
\end{equation}
Again Vieta's formulas imply \eqref{deg:degree-graph}, and
\[
 (\Lambda_1,\Lambda_2,\Lambda_3)=(1,U^{-3},GU^{-2})
\]
recovers \((p,\gamma,s)=(tU,t^2G,t^3)\) exactly.  Uniformly for
\(|G|\le G_{\max}\), all positive units in
\eqref{deg:cubic-degree-dictionary} tend to one and \(\nu\to0\).

The compact-degree positive units have compact closure. Let
\(\tau_d>0\) be their radial bound, and let the positive far-root
bound \(\tau_0\) include its fixed core/high splice reserve. Choose once
\begin{equation}\label{deg:small-ratio-choice}
 0<p_*^2<\tfrac12\min\{\tau_0,\tau_d,
                                  \tau_-\omega_{\min}^2\}.
\end{equation}
Then the degree and positive far-root scales \(\tau=p^2\), and the
negative far-root scale \(\tau=p^2/\omega^2\), lie in their required
ranges. These radial restrictions must be chosen explicitly, since shrinking
the raw coefficients alone does not control the chart ratios.

The chosen chart ranges now cover every coefficient alternative. If \(\gamma\ne0\) and
\(0<k\le2k_\flat\), then:
\begin{itemize}
 \item for \(0<\pi\le\Pi_{\max}\), use
 Propositions~\ref{deg:prop-positive-quadratic} and
 \ref{deg:prop-negative-gamma}, according to the sign of \(\gamma\);
 \item for \(\pi>\Pi_{\rm far}\) and \(\mathfrak q\ge \mathfrak q_\flat\), use
 Proposition~\ref{deg:prop-degree-drop} through
 \eqref{deg:gamma-degree-dictionary};
 \item for \(\pi>\Pi_{\rm far}\), \(\mathfrak q<\mathfrak q_\flat\), and \(\gamma>0\),
 use Proposition~\ref{deg:prop-positive-far};
 \item in the same range with \(\gamma<0\), use
 Proposition~\ref{deg:prop-degree-infinity} through
 \eqref{deg:negative-far-dictionary}.
\end{itemize}
The compact and far \(\pi\)-ranges, and the two \(\mathfrak q\)-ranges, overlap
strictly.  If \(\gamma=0\), or if \(k\ge k_\flat\), then
\[
 |G|^{3/2}=\frac{|\gamma|^{3/2}}{s}\le k_\flat^{-1}.
\]
For \(U\le U_{\rm cub}\), use
Proposition~\ref{deg:prop-compact-cubic}; for \(U\ge U_d\), use
Proposition~\ref{deg:prop-degree-drop} through
\eqref{deg:cubic-degree-dictionary}.  These intervals also overlap.

The remaining radii \(g=|\gamma|\), \(t=s^{1/3}\), and \(s\)
can be restricted independently. This gives eventual \(s_*,\gamma_*>0\);
the common raw cap is fixed before their final reduction and the choice
of \(m\). The closed alternatives include every equality face, and the bands in \(k,\pi,\mathfrak q,U\) overlap strictly. These choices depend only on the coefficients and are made before any
phase point is selected.
\end{proof}

\subsection{Proof of Lemma~\ref{lem:large-ratio-height}}
\label{app:proof-large-ratio-height}

\begin{proof}
At the central section,
\[
 Y_L(0)=p^{-2}Y_0\le K_0(1+p^{-2})\le2K_0.
\]
The graph equation obtained directly from
\eqref{deg:central-field} is
\begin{equation}
 \frac{dY_L}{dX}=
 \frac{Y_L\{q_\gamma X+r_LX^2+
       Cr_L(\gamma+r_LX)Y_L\}}
  {\ell(\gamma+r_LX)Y_L+p^{-2}-X+3\gamma X^2+r_LX^3}.
 \label{eq:large-ratio-graph}
\end{equation}
The denominator is positive on a through graph. At low height it may
approach zero, so we estimate growth only after the height has crossed
a fixed positive threshold.  Choose \(M_0>2K_0+1\) so
that, whenever \(Y_L\ge M_0\), the denominator is at least
\(\ell_0Y_L/2\).  On that region
\[
 Y_L'\le C_{L_c}+Cr_L(|\gamma|+r_LL_c)Y_L
       \le C_{L_c}+Y_L
\]
after a raw coefficient shrink.  Apply this inequality after the first crossing of \(M_0\), or
equivalently to \(\max(M_0,Y_L)\).  Gronwall's
inequality gives
\[
 Y_L(L_c)\le(M_0+C_{L_c}L_c)e^{L_c}=:B_c.
\]
Strict scalar order gives \(Y_L(L_c)>d_L\); since \(d_L\ge0\),
\eqref{deg:large-ratio-gap} follows.  No derivative lower bound for the
map from \(X=0\) to \(X=L_c\) has been used.
\end{proof}

\subsection{Proof of Lemma~\ref{deg:lem-large-source-routing}}
\label{app:proof-large-source-routing}

\begin{proof}
Fix the constant \(K_0\) in Corollary~\ref{rew:cor-high-derivative-signed}.
Its complete derivative bound applies on the connected high-height interval
\begin{equation}
 Y_0>K_0(1+p^2),\qquad 0<s^2Y_0<z_*.
 \label{deg:large-high-height}
\end{equation}
It remains to cover the lower-height complement by the stopped-trace
and strict source intervals.

Use the positive signed trace chart
\begin{equation}
 \alpha _0=-r_L^2U_0d,\qquad
 \alpha _1=r_Lp_L,\qquad
 d=\frac{p_L^2}{U_0p^2},\qquad
 \frac{s^2}{r_L^2}=\frac{p_L^2}{p^2},
 \label{deg:large-trace-dictionary}
\end{equation}
where \(p_L\) and \(U_0\) are positive units.  In the notation of
Proposition~\ref{deg:prop-stopped-trace}, the signed trace parameter is
\[
 a_L^{\rm tr}:=-U_0d=-\frac{p_L^2}{p^2},
 \qquad \alpha _0=r_L^2a_L^{\rm tr}.
\]
The stopped-trace proposition therefore applies with the signed
parameter \(a_L^{\rm tr}\), which includes the positive factor \(d\).
By increasing \(P_*\),
\(a_L^{\rm tr}\) lies in its signed trace block. For an orbit in
the complementary range \(Y_0\le K_0(1+p^2)\),
\[
 Y_L(0)=\frac{s^2Y_0}{r_L^2}
 \le K_0p_L^2(1+p^{-2})\le C.
\]
The height lemma also applies when \(p_L\) ranges over this compact
positive interval: its coordinates are related to the
\(r_L=\alpha_1\) coordinates by
\(X=p_LX_{\rm old}\), \(Y_L=p_L^2Y_{\rm old}\).
Its Gronwall bound is uniform for the corresponding cuts
\(X_{\rm old}=L_c/p_L\) in a fixed compact interval.
Lemma~\ref{lem:large-ratio-height} therefore gives a fixed number
\(B_0\) such that
\begin{equation}
 0<\frac{\mathfrak g_L}{r_L^2}\le B_0.
 \label{deg:large-bounded-gap}
\end{equation}

Take \(\xi_*=\xi_0\) from \eqref{deg:trace-overlap}. The compact annulus
\(\xi_*\le\mathfrak g_L/r_L^2\le B_0\) is covered by finitely many
strict source intervals, overlapping the terminal range
\(0<\mathfrak g_L/r_L^2\le2\xi_*\). On these intervals the coefficients satisfy
\[
 a_L^{\rm tr}=-\frac{p_L^2}{p^2}\in[-P_*^{-2}C,0],
 \qquad p_L\in[1/2,3/2],
\]
and all orbit segments retain strict positive margins. Proposition
\ref{deg:prop-stopped-trace}, with
\eqref{deg:trace-domain-inclusion} and \eqref{deg:trace-overlap}, supplies
their every-point physical cover. Proposition~\ref{sharp:upper-trace}
controls the derivative on the connected union of these intervals. The choices are
\begin{align*}
 &K_0,\ L_c,\ B_c,\ B_0;
   \quad \text{terminal interval},\ \xi_*;\\
 &P_*,\ \text{finite strict cover},\ L_t;
   \quad \text{common raw cap},\ \text{raw coefficient neighbourhood},\ m.
\end{align*}
The strict margins ensure that every later shrink preserves the domains
and overlaps already chosen.
\end{proof}
